\section{De la arquitectura de papers a sistemas reales: BAGEL y \texorpdfstring{$\pi_{0.7}$}{pi0.7}}

El significado de MoT no es solo un paper sobre curvas de pérdida, sino que proporciona un patrón de diseño componible: permite que diferentes modalidades se comuniquen en la misma secuencia, pero asignando parámetros dedicados para generación, comprensión o control. La clase eligió BAGEL y $\pi_{0.7}$ para mostrar cómo este patrón entra en el preentrenamiento multimodal unificado y el control encarnado (embodied), respectivamente. Ambos no son competidores de benchmark para la misma tarea, sino instancias de la misma idea arquitectónica en diferentes fronteras de capacidad.

\subsection{BAGEL: comprensión compartida, generación dedicada}

El capítulo anterior estableció la enrutación consciente de modalidades; esta sección examina su combinación práctica en un modelo multimodal unificado. BAGEL no busca probar que todos los módulos pueden fusionarse, sino demostrar cómo mantener la comprensión (understanding) compartida y preservar una ruta dedicada para la generación.

\lecturefigure{slide-045.jpg}{BAGEL utiliza parámetros conscientes de modalidad para la comprensión y generación unificadas}{Diapositiva 45 del deck oficial; arXiv BAGEL 2505.14683v3; Clase 00:31:39--00:33:33}

\begin{knowledgebox}{Lectura de la figura: BAGEL no es simplemente "un encoder más un head de difusión"}
La izquierda de la diapositiva lista tareas como comprensión, generación y edición; el diagrama del paper enfatiza el preentrenamiento unificado. La explicación clave de la clase se refiere a las rutas de parámetros: la generación de imágenes utiliza parámetros dedicados, mientras que el modelo de lenguaje multimodal del lado de la comprensión sigue procesando conjuntamente entradas de texto e imagen. Esto preserva la razonamiento lingüístico sin exigir que los latentes generativos y las embeddings semánticas sean idénticos.
\end{knowledgebox}

BAGEL también muestra una secuencia importante: el modelo puede planificar en el espacio de texto antes de generar la imagen. El proceso se puede abstraer como
$$
p(v\mid t)=\sum_r p_\theta(r\mid t)\,p_\theta(v\mid t,r),
$$
donde $r$ es una representación de razonamiento/planeación visible o implícita y $v$ es la salida visual. Esta descomposición no implica que todas las imágenes deban preceder a un chain-of-thought escrito, sino que la abstracción lingüística puede servir como interfaz de control y composición.

\begin{importantbox}{Lecciones de diseño de BAGEL}
La comprensión y la generación pueden compartir semántica de alto nivel y contexto de secuencia, mientras utilizan representaciones de imagen y rutas de parámetros diferentes. El valor de un sistema unificado proviene de la componsibilidad entre tareas, no de buscar isomorfismo en cada capa, cada pérdida o cada token.
\end{importantbox}

\teachervoice{00:32:51--00:33:33, el docente describe nuevamente el estado actual como "no haber encontrado una manera efectiva de unificar completamente la generación de imágenes y la comprensión". BAGEL es una división práctica, no significa que este problema abierto haya sido resuelto teóricamente.}

\subsection{\texorpdfstring{$\pi_{0.7}$}{pi0.7}: el contexto multimodal no son solo instrucciones lingüísticas}

Continuando desde la generación de contenido digital, el siguiente paso es permitir que el contexto multimodal influya en acciones reales. Esta sección utiliza $\pi_{0.7}$ para explicar cómo una señal de dirección (steering signal) puede incluir imágenes de subobjetivos y metadatos del episodio, y discute por qué la encarnación requiere diferentes bucles de evaluación cerrados.

\lecturefigure{slide-046.jpg}{$\pi_{0.7}$ impulsa un modelo base consciente de modalidad hacia el control robótico}{Diapositiva 46 del deck oficial; arXiv $\pi_{0.7}$ 2604.15483v2; Clase 00:34:00--00:35:08}

\begin{knowledgebox}{Lectura de la figura: pasar de generar contenido a generar acciones}
La figura muestra múltiples robots, tareas y condiciones de contexto. El steering de $\pi_{0.7}$ no depende solo de una orden lingüística; también puede aprovechar imágenes de subobjetivos, calidad del episodio o información de estrategia. En comparación con la generación de imágenes, la evaluación de las salidas robóticas depende más del éxito de la tarea, la velocidad, las colisiones y la generalización a través de diferentes cuerpos (embodiments); una acción que parece razonable no implica necesariamente un cierre de ejecución exitoso.
\end{knowledgebox}

Si la política condicional en la observación $o_t$, la instrucción lingüística $g$ y el contexto de steering multimodal $c$, entonces
$$
a_t\sim\pi_\theta(a_t\mid o_{\le t},g,c).
$$
$c$ puede codificar "cómo hacer" no solo "qué hacer". Esto indica que el valor de las modalidades no textuales no es solo proporcionar más percepción, sino también convertirse en especificación de estrategia.

\begin{warningbox}{Límite de evidencia de encarnación}
Los resultados robóticos de $\pi_{0.7}$ apoyan la multitarea, la generalización a través de entornos y la steerability, pero no se puede inferir una inteligencia física universal desde la diapositiva. El despliegue real también requiere frecuencia de control, calibración de sensores, recuperación de fallos, restricciones de seguridad y latencia específica del hardware. La clase lo lista como un ejemplo de especialización al estilo MoT, no como una declaración de que el problema robótico está resuelto.
\end{warningbox}

\subsection{Resumen de la sección}

BAGEL y $\pi_{0.7}$ indican que una arquitectura consciente de modalidades puede expandirse en dos direcciones: primero, unificar comprensión, edición y generación de contenido; segundo, convertir el contexto multimodal en condiciones para estrategias de acción. La ley común es\textbf{permitir interacciones de alto nivel y preservar especialización de bajo nivel}. Esto también trae el problema de vuelta a la segunda línea principal de la conferencia: ¿existe realmente un transfer positivo estable entre comprensión, generación y acción?
